\chapter*{前言：不是一条由乱而治的线}
\addcontentsline{toc}{chapter}{前言：不是一条由乱而治的线}

二八〇年春，王濬的舰队抵达建业，孙皓出降。三国以来分立的最高政权重新归于洛阳；三百余年后，隋军入建康，陈后主被俘，政治上的统一又一次完成。这两个场景很容易把本卷写成一条从统一、崩解到复归统一的直线。可是船队、城门和降表只说明谁取得了最高名义。它们不能替我们回答：命令怎样越过河流和山地，粮食、户籍、寺院与军镇怎样被重新接上，流亡家庭又在何处安置。

本卷以八百七十六条事件组成编年脊柱，但不把账本逐条改写为正文。西晋的破裂不是宫廷人物的道德故事；它同关中军费、道路、军镇与继承程序一起失灵。东晋也不是一个把北方原样搬到江南的“小朝廷”：建康是在迁徙、侨置、江东旧有社会和外镇兵力之间不断协商出来的中心。北方所谓“十六国”不是一张静止的国家名单；城郭、牧地、绿洲、河西道路和不同政权的文书网络，使它成为不断重组的世界。

北魏将平城、军镇、编户、佛教工程和洛阳接进一套更大的统治实验，南朝则在州镇、淮汉防线、文书与宗教网络中维系建康。它们的竞争并不只发生在战场。均田、三长、寺院禁断、造像、户籍、屯田和运河，都是国家试图把人和资源变得可见、可调动的办法；政策文本却从来不是地方落实的保证。六镇、侯景之乱和东西魏的分立，正显示资源网络一旦在不同空间获得自己的军事与社会中介，皇帝的诏令便难以重新把它们接合。

书中以《晋书》、南北朝正史和《资治通鉴》维持年序，又让墓志、造像题记、佛道文本、河西和吐鲁番文书、城址与道路遗存提出不同尺度的问题。后出的史书能保存行动的骨架，却不能替代现场档案；一方墓志或一处石窟能校正时间、人物和物质联系，却不能代表全部社会。遇到年代、兵数、疆域、族名、制度执行和人物动机的争议，本卷宁可明确它的证据边界，也不以整齐的故事填补空白。

因此，读者应把二八〇年和五八九年看作两个开口：一个开口让我们看见统一后的接收为何失败，另一个开口让我们看见再统一须经过怎样漫长而不均匀的行政、社会和交通重组。
